\chapter{Opioid Drug Testing}
\section*{What Is This Test?}
Opioid drug testing detects the presence of natural, semi-synthetic, and synthetic opioids and their metabolites in urine, blood, or oral fluid \cite{cone2001alternate}. The standard opiate immunoassay screen (opiates panel) uses antibodies directed against morphine and detects natural opiates: morphine, codeine, and heroin (as 6-monoacetylmorphine/6-MAM, the specific metabolite of heroin, detectable in urine for only 6--12 hours after use). The standard opiate screen does NOT reliably detect semi-synthetic opioids (oxycodone, hydrocodone, oxymorphone, hydromorphone) or synthetic opioids (fentanyl, methadone, buprenorphine, tramadol, tapentadol, meperidine); these require specific immunoassays or confirmatory LC-MS/MS. Opioid metabolism: codeine is metabolized by CYP2D6 to morphine (its active metabolite); heroin is rapidly metabolized to 6-MAM (t1/2 5--30 minutes) then to morphine; fentanyl is metabolized by CYP3A4 to norfentanyl; methadone has a long half-life (15--60 hours) and is excreted as its major metabolite EDDP (2-ethylidene-1,5-dimethyl-3,3-diphenylpyrrolidine) \cite{trescot2008opioid}. Confirmatory testing by gas chromatography-mass spectrometry (GC-MS) or liquid chromatography-tandem mass spectrometry (LC-MS/MS) is essential for definitive identification of specific opioids.
\section*{Alternative Names}
Opiate Screen; Opioid Toxicology Screen; Opiate Urine Test; Opioid Confirmatory Test; Narcotic Screen
\section*{Principle}
Preliminary screening: competitive immunoassay, cutoff 300 ng/mL for opiates. Confirmatory: GC-MS or LC-MS/MS \cite{tietz2012textbook}. Detection windows: morphine/codeine 2--3 days, heroin (6-MAM) 6--12 hours, oxycodone 2--4 days, fentanyl 1--4 days, methadone 2--7 days, buprenorphine 2--5 days \cite{moeller2008urine}.
\section*{History}
Opium has been used medicinally for millennia, and morphine was first isolated from opium by Friedrich Sertürner in the early 1800s. Heroin (diacetylmorphine) was synthesized by Bayer in 1898 and marketed as a cough suppressant before its addictive potential became apparent. The modern era of drug testing began in the 1960s with the development of immunoassays for opiates in urine, and workplace drug testing expanded in the United States after the federal Drug-Free Workplace initiatives of the 1980s. The development of gas chromatography--mass spectrometry (GC-MS) confirmation in the 1970s--1980s established the two-tier (screen and confirm) paradigm still used today. The expansion of chronic opioid prescribing for pain in the 1990s--2000s drove the need for expanded panels covering synthetic and semi-synthetic opioids and for quantitative LC-MS/MS monitoring in pain clinics.
\section*{Why This Test Is Ordered}
\begin{itemize}
  \item Detection of opioid use in the acute care setting, including suspected opioid overdose
  \item Workplace, forensic, and pre-employment drug testing mandated by federal and state programs
  \item Monitoring of patients on prescribed opioids for chronic pain to verify adherence and detect diversion or non-prescribed use
  \item Evaluation of patients in addiction treatment or medication-assisted therapy (methadone, buprenorphine) programs
  \item Perioperative and emergency toxicology workup for altered mental status
  \item Evaluation of suspected neonatal abstinence syndrome in newborns
\end{itemize}
\section*{Primary vs Secondary Test}
The urine opiate immunoassay is a primary screening test for opioid exposure but is not diagnostic of intoxication or impairment. Confirmatory testing by GC-MS or LC-MS/MS is a secondary test required to establish the specific opioid(s) present and to verify a positive screen. Quantitative LC-MS/MS panels are primary tests for monitoring adherence in chronic pain management.
\section*{How This Test Is Performed}
A urine sample is collected under observed or chain-of-custody conditions for forensic and workplace testing. The sample is screened by a competitive immunoassay (point-of-care cup or laboratory analyzer) with a cutoff of 300 ng/mL for opiates. Specimens are checked for validity (temperature, creatinine, pH) to detect adulteration or dilution. Positive screens are confirmed by GC-MS or LC-MS/MS; for pain management monitoring, quantitative LC-MS/MS measures specific opioids and their metabolites.
\section*{What Preparation Is Needed}
None. Patients should inform the collector of prescribed medications, over-the-counter drugs, and recent poppy seed ingestion, which can produce a positive morphine/codeine result. No fasting or timing requirements apply; detection depends on the window since last use (2--3 days for most opioids).
\section*{Contraindications and Precautions}
No contraindications to urine collection. Precautions: (1) a positive opiate screen is not proof of illicit use --- therapeutic opioids, poppy seeds, and some medications (rifampin, quinolones, dextromethorphan) can cause false positives, so every positive screen requires confirmatory testing; (2) a negative screen does not exclude use, because synthetic and semi-synthetic opioids are not detected by the standard opiate assay; (3) samples may be adulterated, diluted, or substituted, so specimen validity testing is essential in workplace settings; (4) results must not be used alone to judge impairment or dosing \cite{wolff1999review}.
\section*{Risks}
No physical risk; the test is noninvasive. Legal, employment, and privacy consequences can follow from results, so chain-of-custody procedures and informed consent are important.
\section*{What Happens After the Test}
A positive screen is confirmed by GC-MS or LC-MS/MS before any definitive interpretation is made. In suspected overdose, naloxone is administered and serum opioid concentrations are not typically required for management. For pain management, a confirmed positive or unexpected result is discussed with the patient and prescribing clinician, and the monitoring plan or treatment may be adjusted.
\section*{Understanding Results}
Positive opiate screen: Must be confirmed by GC-MS/LC-MS/MS. Poppy seed ingestion produces morphine and codeine positive. Heroin use: 6-MAM positive (confirmatory). Fentanyl/methadone/buprenorphine/oxycodone: negative on standard opiate screen despite use. False-positive opiate screen: quinolone antibiotics (levofloxacin, ofloxacin), rifampin, dextromethorphan (rare). Semisynthetic/synthetic opioids require specific immunoassays.
\section*{Normal Results}
Negative for all tested opioids below the cutoff. Opioid detection in urine does not indicate impairment, dose, timing of use, or distinguish therapeutic use from misuse. Pain management monitoring requires LC-MS/MS quantitative opioid/metabolite testing.
\section*{Follow-up Tests}
Confirmatory GC-MS or LC-MS/MS with quantification. Naloxone response in suspected overdose. Serum/plasma opioid concentrations (clinical toxicity).
\section*{References}
\bibliographystyle{unsrtnat}
\bibliography{references}

\clearpage
